\documentclass[11pt]{article}
\usepackage[margin=1in]{geometry}
\usepackage{booktabs}
\usepackage{amsmath}
\usepackage{siunitx}
\usepackage{caption}
\usepackage{hyperref}

\title{Weighting Scheme: Equal Risk Contribution (ERC)}
\author{MetaFVG Mean-Reversion Portfolio -- 14-Instrument Diversified Universe}
\date{\today}

\begin{document}
\maketitle

\section{Scheme Description}

The portfolio blends 14 instruments trading a Spearman-gated MetaFVG mean-reversion
strategy, deliberately assembled across different central banks and currency blocs
(AUDUSD, EURUSD, GBPUSD, NZDUSD, USDCAD, USDCHF, USDJPY, US500, GER40, JP225, HK50,
USDSGD, EURNOK, USDZAR). The current production blend uses equal weights
($w_i = 1/14$ for all $i$). This report tests whether replacing equal weighting with
\emph{Equal Risk Contribution} (ERC), also known as true risk parity, further reduces
realized correlation and improves diversification.

Unlike inverse-volatility weighting, ERC uses the full covariance matrix
$\Sigma = \mathrm{Cov}(r)$ of daily returns, so it accounts for cross-instrument
correlation rather than standalone volatility alone. The goal is to find a weight
vector $w$, with $w_i > 0$ and $\sum_i w_i = 1$, such that every instrument
contributes equally to total portfolio variance.

\subsection{ERC Formulation}

Total portfolio variance is $\sigma_p^2(w) = w^\top \Sigma w$. The marginal
contribution to risk of asset $i$ is
\[
  \mathrm{MRC}_i = \frac{\partial \sigma_p^2}{\partial w_i} = (\Sigma w)_i,
\]
and its risk contribution is
\[
  \mathrm{RC}_i = w_i \cdot \mathrm{MRC}_i = w_i (\Sigma w)_i, \qquad
  \sum_i \mathrm{RC}_i = \sigma_p^2(w).
\]
ERC seeks $w$ such that $\mathrm{RC}_i = \sigma_p^2(w) / n$ for all $i$, i.e.\ every
instrument contributes the same fraction of total variance.

\subsection{Solve Algorithm}

The weights were found with a cyclical multiplicative-update algorithm (no
\texttt{scipy.optimize} dependency, and no BLAS matrix-matrix multiply -- the
$\Sigma w$ product is computed via elementwise broadcast-and-sum rather than \texttt{@}
to avoid a known BLAS/LAPACK crash in this environment):

\begin{align*}
  w^{(0)}_i &= 1/n \\
  \mathrm{MRC}^{(k)} &= \Sigma w^{(k)} \quad \text{(elementwise broadcast + row sum)} \\
  \mathrm{RC}^{(k)}_i &= w^{(k)}_i \cdot \mathrm{MRC}^{(k)}_i \\
  \mathrm{target}^{(k)} &= \frac{1}{n}\sum_i \mathrm{RC}^{(k)}_i \\
  \tilde{w}^{(k+1)}_i &= w^{(k)}_i \sqrt{\mathrm{target}^{(k)} / \mathrm{RC}^{(k)}_i} \\
  w^{(k+1)} &= \tilde{w}^{(k+1)} / \sum_i \tilde{w}^{(k+1)}_i
\end{align*}

iterated up to 2000 times, stopping when $\max_i \mathrm{RC}_i - \min_i \mathrm{RC}_i <
10^{-8} \cdot \sigma_p^2(w)$.

\textbf{Convergence:} the solve converged in \textbf{11 iterations}, well within the
2000-iteration budget. At convergence the spread of risk contributions relative to
target, $(\max_i \mathrm{RC}_i - \min_i \mathrm{RC}_i) / \mathrm{target}$, was
$4.97 \times 10^{-8}$ -- effectively exact equalization of per-instrument risk
contribution.

\section{ERC Weights}

\begin{table}[h]
\centering
\caption{ERC weights, sorted descending. Ratio column is $w_i / (1/14)$.}
\begin{tabular}{lS[table-format=1.5]S[table-format=1.3]}
\toprule
Symbol & {Weight} & {Weight / (1/14)} \\
\midrule
USDSGD & 0.13622 & 1.907 \\
USDCAD & 0.10944 & 1.532 \\
EURUSD & 0.10283 & 1.440 \\
EURNOK & 0.09940 & 1.392 \\
USDCHF & 0.09032 & 1.265 \\
GBPUSD & 0.07949 & 1.113 \\
USDJPY & 0.07403 & 1.036 \\
AUDUSD & 0.05938 & 0.831 \\
NZDUSD & 0.05415 & 0.758 \\
USDZAR & 0.05109 & 0.715 \\
GER40  & 0.04096 & 0.573 \\
US500  & 0.04008 & 0.561 \\
HK50   & 0.03372 & 0.472 \\
JP225  & 0.02888 & 0.404 \\
\midrule
Sum & 1.00000 & 14.000 \\
\bottomrule
\end{tabular}
\end{table}

ERC concentrates weight in the lowest-variance G10 FX pairs (USDSGD, USDCAD, EURUSD,
EURNOK, USDCHF) at roughly 1.1--1.9$\times$ the equal-weight allocation, and pulls
weight away from the highest-variance instruments -- the equity indices (JP225, HK50,
US500, GER40) and the more volatile crosses (NZDUSD, USDZAR, AUDUSD) -- down to
roughly 0.4--0.8$\times$ equal weight.

\section{Correlation Impact}

\begin{table}[h]
\centering
\caption{Realized correlation and diversification impact, equal-weight vs.\ ERC.}
\begin{tabular}{lS[table-format=1.6]S[table-format=1.6]}
\toprule
Metric & {Equal Weight} & {ERC} \\
\midrule
Weighted avg.\ pairwise correlation & 0.015615 & 0.014191 \\
Diversification ratio (DR) & 2.080922 & 1.830624 \\
\bottomrule
\end{tabular}
\end{table}

ERC lowers the weighted average pairwise correlation slightly (0.0156 $\to$ 0.0142,
a $\sim$9\% relative reduction) by tilting away from instruments that co-move more
with the rest of the book. However, the diversification ratio \emph{falls}
(2.081 $\to$ 1.831): DR rewards spreading weight across instruments with high
\emph{standalone} volatility relative to the blended portfolio's volatility, and ERC
does the opposite by design -- it deliberately down-weights the highest-standalone-vol
instruments (equity indices, JP225/HK50/US500/GER40) to equalize their risk
\emph{contribution}. The two diversification metrics move in opposite directions here.

\section{Results}

\begin{table}[h]
\centering
\caption{Quantstats metrics, equal-weight vs.\ ERC blend.}
\begin{tabular}{lS[table-format=1.6]S[table-format=1.6]}
\toprule
Metric & {Equal Weight} & {ERC} \\
\midrule
Sharpe & 0.517941 & 0.311482 \\
Sortino & 0.816977 & 0.487811 \\
Calmar & 0.343017 & 0.191157 \\
CAGR & 0.000183 & 0.000102 \\
Max Drawdown (\%) & -5.325228 & -5.325228 \\
Volatility (\%) & 3.527622 & 3.269655 \\
Tail Ratio & 1.238689 & 1.114092 \\
Ulcer Index & 0.000192 & 0.000183 \\
Kelly Criterion & 0.066763 & 0.045890 \\
VaR (\%) & -0.358268 & -0.334747 \\
CVaR (\%) & -0.715011 & -0.733195 \\
Skew & 1.544839 & 1.948282 \\
Recovery Factor & 2.260123 & 1.259807 \\
\midrule
Win rate & 0.308850 & 0.304636 \\
Best day & 0.000291 & 0.000291 \\
Worst day & -0.000239 & -0.000239 \\
Total return & 0.001204 & 0.000671 \\
\bottomrule
\end{tabular}
\end{table}

\section{Takeaway}

ERC modestly reduces weighted average pairwise correlation (-9\% relative) but this
does not translate into a better risk-adjusted portfolio: Sharpe falls from 0.518 to
0.311, Sortino from 0.817 to 0.488, and total return drops by more than 40\%, while
maximum drawdown is unchanged. The diversification ratio actually \emph{worsens}
(2.081 $\to$ 1.831), because ERC's risk-contribution-equalizing objective directly
conflicts with DR's standalone-volatility-weighted definition of diversification in
this universe. On this data, equal weighting is the better blend -- ERC trades away
return concentrated in the equity-index legs (which it underweights) for a small,
not-obviously-worthwhile correlation reduction.

\end{document}
